\subsubsection{Flat modules and flat ring maps}
\label{editorial-source-section-flat}

\noindent
One often used result is that if $M = \colim_{i\in \mathcal{I}} M_i$
is a colimit of $R$-modules and if $N$ is an $R$-module then
$$
M \otimes N
=
\colim_{i\in \mathcal{I}} M_i \otimes_R N,
$$
see Lemma \ref{editorial-source-lemma-tensor-products-commute-with-limits}.
This property is usually expressed by saying
that {\it $\otimes$ commutes with colimits}.
Another often used result is that if $0 \to N_1 \to N_2 \to N_3 \to 0$
is an exact sequence and if $M$ is any $R$-module, then
$$
M \otimes_R N_1
\to
M \otimes_R N_2
\to
M \otimes_R N_3
\to
0
$$
is still exact, see Lemma \ref{algebra-lemma-tensor-product-exact}.
Both of these properties tell us that the functor
$N \mapsto M \otimes_R N$ {\it is right exact}.
See Categories, Section \ref{categories-section-exact-functor}
and Homology, Section \ref{homology-section-functors}.
An $R$-module $M$ is flat if $N \mapsto N \otimes_R M$ is also left exact,
i.e., if it is exact. Here is the precise definition.

\begin{definition}
\label{editorial-source-definition-flat}
Let $R$ be a ring.
\begin{enumerate}
\item An $R$-module $M$ is called {\it flat} if whenever
$N_1 \to N_2 \to N_3$ is an exact sequence of $R$-modules
the sequence $M \otimes_R N_1 \to M \otimes_R N_2 \to M \otimes_R N_3$
is exact as well.
\item An $R$-module $M$ is called {\it faithfully flat} if the
complex of $R$-modules
$N_1 \to N_2 \to N_3$ is exact if and only if
the sequence $M \otimes_R N_1 \to M \otimes_R N_2 \to M \otimes_R N_3$
is exact.
\item A ring map $R \to S$ is called {\it flat} if
$S$ is flat as an $R$-module.
\item A ring map $R \to S$ is called {\it faithfully flat} if
$S$ is faithfully flat as an $R$-module.
\end{enumerate}
\end{definition}

\noindent
Here is an example of how you can use the flatness condition.

\begin{lemma}
\label{editorial-source-lemma-flat-intersect-ideals}
Let $R$ be a ring. Let $I, J \subset R$ be ideals. Let $M$ be a flat
$R$-module. Then $IM\cap JM=(I\cap J)M$.
More generally, if $A_1,\ldots,A_r$ are submodules of an
$R$-module $N$, their tensors with $M$ identify with
submodules of $N\otimes_R M$, and
$$
\left(\bigcap_{i=1}^r A_i\right)\otimes_R M
\ \cong\ \bigcap_{i=1}^r(A_i\otimes_R M)
$$
under these original inclusions, for every finite $r\geq1$.
\end{lemma}

\begin{proof}
In the original ideal case use the exact sequence
$$
0\longrightarrow I\cap J\longrightarrow R
\longrightarrow R/I\oplus R/J,
$$
whose last map is $r\mapsto(r+I,r+J)$. Tensoring with $M$
preserves exactness. The maps
$$
R\otimes_R M\longrightarrow M,\quad r\otimes m\longmapsto rm,
\qquad
(R/I)\otimes_R M\longrightarrow M/IM,\quad
(r+I)\otimes m\longmapsto rm+IM
$$
are isomorphisms with inverses $m\mapsto1\otimes m$ and
$m+IM\mapsto(1+I)\otimes m$ respectively. The latter is
well-defined since each $a m$ with $a\in I$ maps to
$(a+I)\otimes m=0$; their composites fix the original tensor
generators and residue classes. The same formula applies to
$J$. Tensor product respects the indicated finite direct sum
by the component projections and inclusions. Thus the exact
sequence becomes
$$
0\longrightarrow (I\cap J)\otimes_R M\longrightarrow M
\longrightarrow M/IM\oplus M/JM.
$$
The middle kernel is $IM\cap JM$. The first map sends
$a\otimes m$ to $am$, so its image is exactly $(I\cap J)M$.
This proves the original equality.

For the general statement, tensoring $0\to A_i\to N$ with
$M$ proves that each original map $A_i\otimes_R M\to N\otimes_R M$
is injective. For two submodules $A,B\subseteq N$, use
$$
0\longrightarrow A\cap B\longrightarrow N
\longrightarrow (N/A)\oplus(N/B),\qquad
n\longmapsto(n+A,n+B).
$$
After tensoring, exactness identifies
$(A\cap B)\otimes_R M$ with the kernel of the two quotient
maps. Right exactness for $A\to N\to N/A\to0$ identifies
the kernel of the first quotient map with the image of
$A\otimes_R M$; the sequence for $B$ identifies the other
kernel. The kernel of the pair is their intersection, proving
the displayed isomorphism for two submodules. Its map is
induced by their original inclusions, so induction, intersecting
the first $r-1$ submodules with $A_r$, proves the finite case.
For $r=1$ it is the identity under the given inclusion.
\end{proof}

\begin{lemma}
\label{editorial-source-lemma-colimit-flat}
Let $R$ be a ring. Let $\{M_i, \varphi_{ii'}\}$ be a directed system of
flat $R$-modules. Then $\colim_i M_i$ is a flat $R$-module.
\end{lemma}

\begin{proof}
Write $M=\colim_i M_i$ with its original structure maps
$\lambda_i:M_i\to M$. For each module $N$, the natural
isomorphism
$$
\colim_i(N\otimes_R M_i)\longrightarrow N\otimes_R M,
\qquad [n\otimes m_i]\longmapsto n\otimes\lambda_i(m_i)
$$
is Lemma \ref{editorial-source-lemma-tensor-products-commute-with-limits};
it commutes with every map of $N$ because both composites
send the displayed generator to the tensor of its original
image. Given an exact complex $N_1\to N_2\to N_3$, flatness
of each $M_i$ makes its tensor complex exact at every stage.
Lemma \ref{algebra-lemma-directed-colimit-exact} gives exactness of
their directed colimit, and the displayed natural isomorphisms
identify it with the tensor complex for $M$. This proves
flatness without assuming that any transition map is injective.
\end{proof}

\begin{example}
\label{editorial-source-example-flat-infinite-intersections}
The finite-intersection assertion above cannot be extended to
arbitrary families of ideals for all flat modules. Take
$R=\mathbf Z$, $M=\mathbf Q$ and $I_n=2^n\mathbf Z$ for
$n\geq1$. The module $\mathbf Q$ is flat by
Lemma \ref{editorial-source-lemma-colimit-flat}: it is the colimit of copies
of $\mathbf Z$ indexed by $n\geq1$, with transition from
$n$ to $n+1$ multiplication by $n+1$ and structure map
$a\mapsto a/n!$. These maps are compatible. Every rational
$a/b$, $b\ne0$, has such a representation since $|b|$ divides
$n!$ for $n\geq|b|$. Equality $a/n!=a'/m!$ is equivalent
after a common stage $k\geq n,m$ to
$a(k!/n!)=a'(k!/m!)$, which is exactly the colimit relation.
Each copy of $\mathbf Z$ is flat because its tensor functor
is the identity via $a\otimes n\mapsto an$ and inverse
$n\mapsto1\otimes n$.

An integer divisible by every $2^n$ must be zero: for a
nonzero integer $a$, choose $n$ with $2^n>|a|$; it cannot
be a nonzero multiple of $2^n$. Thus $\bigcap_{n\geq1}I_n=0$.
On the other hand $I_nM=\mathbf Q$, since every original
$q\in\mathbf Q$ equals $2^n(q/2^n)$. Therefore
$$
\left(\bigcap_{n\geq1}I_n\right)M=0
\quad\hbox{whereas}\quad
\bigcap_{n\geq1}(I_nM)=\mathbf Q.
$$
\end{example}

\begin{lemma}
\label{editorial-source-lemma-composition-flat}
A composition of (faithfully) flat ring maps is
(faithfully) flat.
If $R \to R'$ is (faithfully) flat, and $M'$ is a
(faithfully) flat $R'$-module, then $M'$ is a
(faithfully) flat $R$-module.
\end{lemma}

\begin{proof}
For every original $R$-module $N$ there are mutually inverse
maps
$$
\begin{aligned}
\theta_N:M'\otimes_{R'}(R'\otimes_R N)&\longrightarrow M'\otimes_R N,
&m'\otimes(r'\otimes n)&\longmapsto r'm'\otimes n,\\
\eta_N:M'\otimes_R N&\longrightarrow M'\otimes_{R'}(R'\otimes_R N),
&m'\otimes n&\longmapsto m'\otimes(1\otimes n).
\end{aligned}
$$
The first respects $R'$-balancing because multiplication by
an $R'$-scalar can move between $m'$ and $r'$. It respects
$R$-balancing in the inner tensor by the restricted $R$-action
on $M'$. For the inverse, the equality
$(rm')\otimes(1\otimes n)=m'\otimes(\varphi(r)\otimes n)
=m'\otimes(1\otimes rn)$ proves $R$-balancing.
The composite $\theta_N\eta_N$ fixes $m'\otimes n$, and
$\eta_N\theta_N$ sends a generator to
$r'm'\otimes(1\otimes n)=m'\otimes(r'\otimes n)$.
Both maps commute with each homomorphism of $N$ on those
generators, so they are natural. Here $\varphi:R\to R'$
is the original ring map.

If $R'$ is flat over $R$ and $M'$ flat over $R'$, apply
their tensor functors successively to an exact complex
$N_1\to N_2\to N_3$. Both steps preserve exactness, and
$\theta_{N_i}$ identifies the resulting complex with
$M'\otimes_R N_1\to M'\otimes_R N_2\to M'\otimes_R N_3$.
This proves flatness over $R$. If both hypotheses are faithful,
an arbitrary original complex is exact if and only if the
first tensor complex is exact, and this in turn holds if
and only if the second tensor complex is exact. The same
natural maps prove faithful flatness over $R$. Taking $M'$
to be the target ring of a second ring map gives the first
assertion, for flat and faithfully flat maps alike.
\end{proof}

\begin{lemma}
\label{editorial-source-lemma-flat}
Let $M$ be an $R$-module. The following are equivalent:
\begin{enumerate}
\item
\label{editorial-source-item-flat}
$M$ is flat over $R$.
\item
\label{editorial-source-item-injective}
for every injection of $R$-modules $N \subset N'$
the map $N \otimes_R M \to N'\otimes_R M$ is injective.
\item
\label{editorial-source-item-f-ideal}
for every ideal $I \subset R$ the map
$I \otimes_R M \to R \otimes_R M = M$ is injective.
\item
\label{editorial-source-item-ffg-ideal}
for every finitely generated ideal $I \subset R$
the map $I \otimes_R M \to R \otimes_R M = M$ is injective.
\end{enumerate}
\end{lemma}

\begin{proof}
The successive implications
$(\ref{editorial-source-item-flat})\Rightarrow(\ref{editorial-source-item-injective})
\Rightarrow(\ref{editorial-source-item-f-ideal})\Rightarrow(\ref{editorial-source-item-ffg-ideal})$
follow directly: apply flatness to $0\to N\to N'$, then
specialize to inclusions $I\subseteq R$, and then to finitely
generated ideals. We prove
$(\ref{editorial-source-item-ffg-ideal})\Rightarrow(\ref{editorial-source-item-flat})$.
Suppose that $N_1 \to N_2 \to N_3$ is exact.
Let $K = \Ker(N_2 \to N_3)$ and $Q = \Im(N_2 \to N_3)$.
Then we get maps
$$
N_1 \otimes_R M \to
K \otimes_R M \to
N_2 \otimes_R M \to
Q \otimes_R M \to
N_3 \otimes_R M
$$
Observe that the first and third arrows are surjective. Thus if we show
that the second and fourth arrows are injective, then we are
done\footnote{Here is the argument in more detail:
Assume that we know that the second and fourth arrows are
injective. Lemma \ref{algebra-lemma-tensor-product-exact} (applied
to the exact sequence $K \to N_2 \to Q \to 0$) yields that
the sequence $K \otimes_R M \to N_2 \otimes_R M \to
Q \otimes_R M \to 0$ is exact. Hence,
$\Ker \left(N_2 \otimes_R M \to Q \otimes_R M\right)
= \Im \left(K \otimes_R M \to N_2 \otimes_R M\right)$.
Since
$\Im \left(K \otimes_R M \to N_2 \otimes_R M\right)
= \Im \left(N_1 \otimes_R M \to N_2 \otimes_R M\right)$
(due to the surjectivity of $N_1 \otimes_R M \to
K \otimes_R M$) and
$\Ker \left(N_2 \otimes_R M \to Q \otimes_R M\right)
= \Ker \left(N_2 \otimes_R M \to N_3 \otimes_R M\right)$
(due to the injectivity of $Q \otimes_R M \to
N_3 \otimes_R M$), this becomes
$\Ker \left(N_2 \otimes_R M \to N_3 \otimes_R M\right)
= \Im \left(N_1 \otimes_R M \to N_2 \otimes_R M\right)$,
which shows that the functor $- \otimes_R M$ is exact,
whence $M$ is flat.}.
Hence it suffices to show that $- \otimes_R M$ transforms
injective $R$-module maps into injective $R$-module maps.

\medskip\noindent
Assume $K \to N$ is an injective $R$-module map and
let $x \in \Ker(K \otimes_R M \to N \otimes_R M)$.
We have to show that $x$ is zero.
The $R$-module $K$ is the union of its finite
$R$-submodules; hence, $K \otimes_R M$ is
the colimit of $R$-modules of the form
$K_i \otimes_R M$ where $K_i$ runs over all finite
$R$-submodules of $K$
(because tensor product commutes with colimits).
Thus, for some $i$ our $x$ comes from an element
$x_i \in K_i \otimes_R M$. Thus we may assume that $K$
is a finite $R$-module. Assume this. We regard the
injection $K \to N$ as an inclusion, so that
$K \subset N$.

\medskip\noindent
The $R$-module $N$ is the union of its finite
$R$-submodules that contain $K$. Hence, $N \otimes_R M$
is the colimit of $R$-modules of the form
$N_i \otimes_R M$ where $N_i$ runs over all finite
$R$-submodules of $N$ that contain $K$
(again since tensor product commutes with colimits).
Notice that this is a colimit over a directed system
(since the sum of two finite submodules of $N$ is
again finite).
Hence, (by Lemma \ref{algebra-lemma-zero-directed-limit})
the element $x \in K \otimes_R M$ maps to
zero in at least one of these $R$-modules
$N_i \otimes_R M$ (since $x$ maps to zero
in $N \otimes_R M$).
Thus we may assume $N$ is a finite $R$-module.

\medskip\noindent
Assume $N$ is a finite $R$-module. Choose a surjection
$R^{\oplus n}\to N$, put $L$ equal to its kernel, and let
$L'$ be the inverse image of the original $K\subseteq N$.
Then $L\subseteq L'\subseteq R^{\oplus n}$,
$N=R^{\oplus n}/L$ and $K=L'/L$, under the quotient maps
just specified. For any submodule $G\subseteq R^{\oplus n}$,
the map $G\otimes_R M\to M^{\oplus n}$ sends
$(g_1,\ldots,g_n)\otimes m$ to $(g_1m,\ldots,g_nm)$.
Suppose these maps are injective for $G=L,L'$, and denote
them by $\beta$ and $\beta'$ respectively. Tensoring the
original inclusion gives the map in the direction
$$
\alpha:L\otimes_R M\longrightarrow L'\otimes_R M,
\qquad \ell\otimes m\longmapsto\ell\otimes m,
\qquad \beta'\alpha=\beta.
$$
It is injective since $\beta$ is injective. Right exactness
identifies the original $K\otimes_R M$ with
$(L'\otimes_R M)/\alpha(L\otimes_R M)$ and $N\otimes_R M$
with $M^{\oplus n}/\beta(L\otimes_R M)$.
The induced quotient map is injective: if
$v\in L'\otimes_R M$ has $\beta'(v)=\beta(u)$ for some
$u\in L\otimes_R M$, then
$\beta'(v-\alpha(u))=0$, hence $v=\alpha(u)$ by injectivity
of $\beta'$. This proves the required injection using the
original inclusions, without a reverse map from $L'$ to $L$.

\medskip\noindent
Thus it suffices to show that $L \otimes_R M \to M^{\oplus n}$
is injective when $L \subset R^{\oplus n}$ is an $R$-submodule.
We do this by induction on $n$. If $n=0$, both $L$ and
$M^{\oplus n}$ are zero, so the assertion holds. The base
case $n=1$ is proved below.
For the induction step assume $n > 1$ and set
$L'=L\cap(R\oplus0^{\oplus n-1})$. It is isomorphic to
the ideal $\{a\in R:(a,0,\ldots,0)\in L\}$ by the first
coordinate map, with inverse $a\mapsto(a,0,\ldots,0)$.
The last $n-1$ coordinate projection on $L$ has kernel $L'$,
so it induces the injection
$L''=L/L'\to R^{\oplus n-1}$ sending the class of
$(a_1,\ldots,a_n)$ to $(a_2,\ldots,a_n)$.
It is well-defined since elements of $L'$ have all these
coordinates zero, and injective by the kernel equality.
We obtain a diagram
$$
\xymatrix{
&
L' \otimes_R M \ar[r] \ar[d] &
L \otimes_R M \ar[r] \ar[d] &
L'' \otimes_R M \ar[r] \ar[d] &
0 \\
0 \ar[r] &
M \ar[r] &
M^{\oplus n} \ar[r] &
M^{\oplus n - 1} \ar[r] & 0
}
$$
The left vertical map is induced by the first coordinate,
the right by the last $n-1$ coordinates, and the lower maps
are the corresponding inclusion and projection. The upper
row is right exact, the lower row is exact, and the squares
commute on every tensor generator. The left and right
vertical maps are injective by the base case and induction.
If $v\in L\otimes_R M$ maps to zero in $M^{\oplus n}$,
its image in $L''\otimes_R M$ maps to zero on the right and
is therefore zero. Lift $v$ to $u\in L'\otimes_R M$ using
upper-row exactness. Its image in $M$ maps to zero in
$M^{\oplus n}$ and hence is zero, since the lower left map
is injective. Left vertical injectivity now gives $u=0$ and
then $v=0$, proving the middle injectivity.

\medskip\noindent
The base case of the induction above is when $L \subset R$ is an ideal.
In other words, we have to show that $I \otimes_R M \to M$ is injective
for any ideal $I$ of $R$. We know this is true when $I$ is finitely
generated. However, $I = \bigcup I_\alpha$ is the union of the
finitely generated ideals $I_\alpha$ contained in it. In other words,
$I=\colim_\alpha I_\alpha$, over the directed set of finitely
generated subideals; the sum of two such subideals contains
both. Tensor-colimit compatibility gives
$I\otimes_R M=\colim_\alpha(I_\alpha\otimes_R M)$.
If an element of this colimit maps to zero in $M$, choose
a representative $v_\alpha\in I_\alpha\otimes_R M$.
Its image in $M$ is already zero. The assumed injectivity
at that stage gives $v_\alpha=0$, so the original colimit
element is zero. This proves the base case for every ideal,
and completes the stated criterion.
\end{proof}

\begin{lemma}
\label{editorial-source-lemma-colimit-rings-flat}
Let $\{R_i, \varphi_{ii'}\}$ be a system of rings over the directed set $I$.
Let $R = \colim_i R_i$.
\begin{enumerate}
\item If $M$ is an $R$-module such that $M$ is flat as an $R_i$-module
for all $i$, then $M$ is flat as an $R$-module.
\item For $i \in I$ let $M_i$ be a flat $R_i$-module and
for $i' \geq i$ let $f_{ii'} : M_i \to M_{i'}$ be a $\varphi_{ii'}$-linear
map such that $f_{i' i''} \circ f_{i i'} = f_{i i''}$. Then
$M = \colim_{i \in I} M_i$ is a flat $R$-module.
\end{enumerate}
\end{lemma}

\begin{proof}
Write $\lambda_i:R_i\to R$ for the original ring maps and
$\mu_i:M_i\to M$ for the module structure maps. Part (1)
is the instance of part (2) with every $M_i=M$, acted on
through $\lambda_i$, and all underlying transition maps the
identity. They are semilinear because
$\lambda_{i'}\varphi_{ii'}=\lambda_i$.

We first specify the colimit constructions in part (2).
Only a nonempty directed indexing set is needed; none of
the arguments uses antisymmetry of its order. For a system
of abelian groups $B_i$ with transition maps $b_{ij}$,
the colimit is the direct sum of the $B_i$ modulo the
subgroup generated by the original relations
$\iota_i(x)-\iota_j(b_{ij}(x))$ for $i\leq j$.
Every element has a representative at a single stage:
move its finitely many summands to a common upper stage
and add their images there. Two representatives $x_i,x_j$
are equal in the colimit exactly when their images agree
at some common upper stage. To prove the nontrivial
direction, express their difference as a finite sum of
the generating relations and move every stage occurring
in that finite expression to one upper stage; compatibility
kills each relation there. The converse follows from the
defining transition relations. In particular a stage
element represents zero exactly when its image is zero
at a later stage. This description permits noninjective
transition maps throughout.

For $r=\lambda_i(r_i)$ and $m=\mu_j(m_j)$ choose $k\geq i,j$
and define the original scalar action by
$$
rm=\mu_k\bigl(\varphi_{ik}(r_i)f_{jk}(m_j)\bigr).
$$
Changing either representative or the chosen upper stage
does not change this element: move the relevant equalities
to one further stage and use semilinearity and compatibility.
The distributive laws and associativity of the action follow
by moving all finitely many entries of each identity to one
stage and applying its module axiom. The unit acts as the
identity because $[f_{ik}(m_i)]=[m_i]$.

Here is the varying-base tensor comparison needed below.
Let $A_i$ be $R_i$-modules with compatible semilinear maps
$h_{ij}:A_i\to A_j$, and put $A=\colim_i A_i$ with the
same scalar construction and structure maps $\alpha_i$.
There are transition maps
$$
\tau_{ij}:A_i\otimes_{R_i}M_i\longrightarrow A_j\otimes_{R_j}M_j,
\qquad a_i\otimes m_i\longmapsto h_{ij}(a_i)\otimes f_{ij}(m_i).
$$
They are well-defined: the two images of the balancing
relation $r_ia_i\otimes m_i=a_i\otimes r_im_i$ agree by
$$
\varphi_{ij}(r_i)h_{ij}(a_i)\otimes f_{ij}(m_i)
=h_{ij}(a_i)\otimes\varphi_{ij}(r_i)f_{ij}(m_i).
$$
Let $T=\colim_i(A_i\otimes_{R_i}M_i)$ with maps $\theta_i$.
The first comparison is
$$
F:T\longrightarrow A\otimes_R M,\qquad
\theta_i\left(\sum_{\ell=1}^q a_{i,\ell}\otimes m_{i,\ell}\right)
\longmapsto\sum_{\ell=1}^q
\alpha_i(a_{i,\ell})\otimes\mu_i(m_{i,\ell}).
$$
Each stage formula respects balancing and every transition
relation, so it defines a map of abelian groups. The scalar
formula above makes it $R$-linear by moving a scalar and
the tensor representative to one common stage.
For the inverse, define an additive balanced pairing on
representatives by
$$
B\bigl(\alpha_i(a_i),\mu_j(m_j)\bigr)
=\theta_k\bigl(h_{ik}(a_i)\otimes f_{jk}(m_j)\bigr),
\qquad k\geq i,j.
$$
Its value is independent of representatives and $k$:
move any two eventual equalities of representatives and
both choices of $k$ to one further stage, where the tensors
are identical. Additivity is checked at a common stage for
the two summands. For $r=\lambda_\ell(r_\ell)$, move
$r,a,m$ to one stage $k$. Then
$$
B(ra,m)=\theta_k\bigl(\varphi_{\ell k}(r_\ell)a_k\otimes m_k\bigr)
=\theta_k\bigl(a_k\otimes\varphi_{\ell k}(r_\ell)m_k\bigr)
=B(a,rm).
$$
Thus $B$ induces an $R$-linear map $G:A\otimes_R M\to T$.
On a pure tensor $FG$ gives
$\alpha_k(h_{ik}(a_i))\otimes\mu_k(f_{jk}(m_j))
=\alpha_i(a_i)\otimes\mu_j(m_j)$.
On a stage tensor $GF$ gives its image at a common later
stage, which is the same class $\theta_i(a_i\otimes m_i)$.
Finite sums of these generators prove both composites are
the identity. Hence the exact comparison, with both maps
specified, is
$$
\colim_i(A_i\otimes_{R_i}M_i)
\ \cong\ (\colim_i A_i)\otimes_R(\colim_i M_i).
$$

For completeness, stagewise injections induce an injection
on these directed colimits. Indeed, for a compatible family
of injections $u_i:X_i\to Y_i$, if $[x_i]$ maps to zero,
then at a later stage $j$ compatibility gives $u_j(x_j)=0$.
Stagewise injectivity gives $x_j=0$, hence $[x_i]=0$.
This uses no injectivity of the transition maps.

Now retain the original finitely generated ideal
$\mathfrak a=(a_1,\ldots,a_n)\subseteq R$. Choose one stage
$i$ containing representatives $a_{\ell,i}$ of all its
generators, and put
$$
\mathfrak a'=(a_{1,i},\ldots,a_{n,i})\subseteq R_i,
\qquad
\mathfrak a_j=\mathfrak a'R_j
=(\varphi_{ij}(a_{1,i}),\ldots,\varphi_{ij}(a_{n,i}))
\quad(j\geq i).
$$
The tail $j\geq i$ is cofinal: every original representative,
and every eventual equality, can be moved to a stage above
both its original stage and $i$. It has the same ring and
module colimits. The map
$$
\delta:\colim_{j\geq i}\mathfrak a_j\longrightarrow\mathfrak a,
\qquad [x_j]\longmapsto\lambda_j(x_j)
$$
is surjective. For $x=\sum_{\ell=1}^n r_\ell a_\ell$,
choose a stage $j\geq i$ representing all $r_\ell$ and use
$x_j=\sum_{\ell=1}^n r_{\ell,j}\varphi_{ij}(a_{\ell,i})$.
It is injective: if $\lambda_j(x_j)=0$ in $R$, some
$k\geq j$ has $\varphi_{jk}(x_j)=0$ in $R_k$; because
$\mathfrak a_k$ is an actual ideal in that ring, the image
is zero in $\mathfrak a_k$ itself. This proves that the
inverse formula $x\mapsto[x_j]$ is independent of its
expression and representatives; equal expressions differ
by an element eventually zero in the original ring system.
Both maps respect the scalar action and fix the represented
elements in their composites.

Apply the proved tensor comparison to $A_j=\mathfrak a_j$.
It identifies
$$
\colim_{j\geq i}(\mathfrak a_j\otimes_{R_j}M_j)
\ \cong\ \mathfrak a\otimes_R M.
$$
For every $j$, flatness of $M_j$ makes
$\mathfrak a_j\otimes_{R_j}M_j\to M_j$ injective. Its
formula is $x_j\otimes m_j\mapsto x_jm_j$, and compatibility
of these maps is exactly semilinearity of $f_{jk}$.
Their colimit is injective by the element argument above.
Under the displayed comparison it is the original map
$\mathfrak a\otimes_R M\to M$, $a\otimes m\mapsto am$.
Lemma \ref{editorial-source-lemma-flat} now proves flatness of $M$.
The zero ideal, including a presentation with no generators,
has zero tensor product and requires no exception.

The comparison also proves the full injection criterion
directly. For an injection $N'\to N$ of $R$-modules, regard
both as constant underlying systems with restricted
$R_i$-actions and identity transition maps. Each
$N'\otimes_{R_i}M_i\to N\otimes_{R_i}M_i$ is injective.
Their colimit comparison is exactly
$N'\otimes_R M\to N\otimes_R M$, hence it is injective.
This gives a second proof with all original modules retained.

The explicitly displayed composition law for $f_{ij}$ also
suffices if identities on the diagonal were not stipulated.
Indeed $e_i=f_{ii}$ is an idempotent $R_i$-linear map and
$$
M_i=\operatorname{Im}(e_i)\oplus\operatorname{Ker}(e_i),
\qquad m_i=e_i(m_i)+(m_i-e_i(m_i)).
$$
The intersection of the two summands is zero since $e_i$
is the identity on its image. Each summand of a flat module
is flat: its tensor map on any injection is a direct summand
of the injective tensor map for $M_i$. Compatibility gives
$f_{ij}=e_jf_{ij}=f_{ij}e_i$, so the restrictions to
$\operatorname{Im}(e_i)$ form a genuine directed system.
The colimit by all the original relations is isomorphic to
the colimit of these images. The maps send $[m_i]$ to
$[e_i(m_i)]$ and an image representative to $[m_i]$ in the
original colimit; they respect transitions by the displayed
identities and are inverse because $[e_i(m_i)]=[m_i]$ is
one of the original relations. Thus no extra identity
hypothesis is needed for the flatness conclusion.

\medskip\noindent
We record exactly what is additionally needed for faithful
flatness. For a maximal ideal $\mathfrak m\subset R$ set
$\mathfrak m_i=\lambda_i^{-1}(\mathfrak m)$. Each is a
proper prime, though it need not be maximal in $R_i$, and
$\varphi_{ij}^{-1}(\mathfrak m_j)=\mathfrak m_i$.
Consequently the induced maps $R_i/\mathfrak m_i\to R_j/\mathfrak m_j$
are injective. The comparison
$$
\colim_i(R_i/\mathfrak m_i)\longrightarrow R/\mathfrak m,
\qquad [r_i+\mathfrak m_i]\longmapsto\lambda_i(r_i)+\mathfrak m
$$
is surjective by representatives and injective because an
element mapping to zero already has $r_i\in\mathfrak m_i$.
Its inverse sends any representative $\lambda_i(r_i)+\mathfrak m$
to $[r_i+\mathfrak m_i]$; the same zero test proves independence
of the choice. Applying the varying-base tensor comparison
gives
$$
M/\mathfrak mM\ \cong\
\colim_i(M_i/\mathfrak m_iM_i).
$$
Here the stage quotient-tensor isomorphism sends
$(r_i+\mathfrak m_i)\otimes u_i$ to
$r_iu_i+\mathfrak m_iM_i$ and has inverse
$u_i+\mathfrak m_iM_i\mapsto(1+\mathfrak m_i)\otimes u_i$.
The same formulas identify the receiving quotient over $R$.
Balancing makes both maps well-defined, and both composites
fix the indicated generators. Thus the original transition
on the right is
$u_i+\mathfrak m_iM_i\mapsto f_{ij}(u_i)+\mathfrak m_jM_j$,
and the receiving map is
$u_i+\mathfrak m_iM_i\mapsto\mu_i(u_i)+\mathfrak mM$.
The eventual-zero rule now proves
$$
\mu_i(u_i)+\mathfrak mM=0
\quad\Longleftrightarrow\quad
f_{ij}(u_i)\in\mathfrak m_jM_j\text{ for some }j\geq i.
$$

A flat module is faithfully flat exactly when all these
maximal-ideal quotients are nonzero. Here is the full
argument, using the definition by exactness reflection.
Faithful flatness detects nonzero modules by applying that
definition to $0\to N\to0$, and in particular detects
$N=R/\mathfrak m$. Conversely suppose all the quotients
are nonzero. Given $0\ne n\in N$, let
$J=\operatorname{Ann}_R(n)$, choose a maximal ideal
$\mathfrak m\supseteq J$, and retain the injection
$R/J\to N$, $r+J\mapsto rn$. The module $M/JM$ surjects
onto the nonzero $M/\mathfrak mM$, so it is nonzero.
Flatness makes $(R/J)\otimes_R M=M/JM\to N\otimes_R M$
injective, proving $N\otimes_R M\ne0$ for every nonzero $N$.

To deduce exactness reflection, take an original complex
$N_1\xrightarrow{d_1}N_2\xrightarrow{d_2}N_3$ and write
$J=\operatorname{Im}(d_1)\subseteq K=\operatorname{Ker}(d_2)$.
Flatness embeds $J\otimes_R M$ and $K\otimes_R M$ in
$N_2\otimes_R M$. The exact sequences defining image and
kernel identify them with the tensor image and tensor
kernel, respectively: the map $N_1\to J$ is surjective,
and $0\to K\to N_2\to N_3$ stays exact. Right exactness
for $J\to K\to K/J\to0$ then identifies
$$
(K/J)\otimes_R M\ \cong\
\operatorname{Ker}(d_2\otimes\mathrm{id}_M)\big/
\operatorname{Im}(d_1\otimes\mathrm{id}_M).
$$
These are the maps induced by the original inclusions and
quotient $K\to K/J$, so the identification respects those
original submodules. If the tensor complex is exact, this
quotient is zero. Nonzero-module detection gives $K/J=0$,
hence the original complex is exact. This proves the criterion.

Combining it with the exact quotient comparison proves:
$M$ is faithfully flat over $R$ if and only if, for every
maximal $\mathfrak m\subset R$, there are an index $i$ and
$u_i\in M_i$ such that
$$
f_{ij}(u_i)\notin\mathfrak m_jM_j
\quad\hbox{for every }j\geq i.
$$
This quantifier keeps the same original element at every
later stage. It cannot be replaced by merely having a
nonzero quotient at each stage.

For example take constant $R_i=k$ and $M_i=k$ over a
nonzero field, indexed by $i\in\mathbf Z_{\geq0}$, with
identity ring maps and $f_{ij}=0$ for $i<j$, $f_{ii}$ the
identity. Each stage module is faithfully flat, but every
element dies at the next stage, so $R=k$ and $M=0$.
Even injective module transitions do not suffice. Take
$R_i=M_i=\mathbf Z$ with identity ring maps and
$f_{ij}(a)=2^{j-i}a$. The map $[a\text{ at }i]\mapsto a/2^i$
identifies $M$ with $\mathbf Z[1/2]$: it is surjective by
the description of its elements, and equality of two
fractions is exactly equality after multiplication to a
common power-of-two denominator. Every stage is faithfully
flat and every transition injective, but $M/2M=0$ since
$a/2^i=2(a/2^{i+1})$, so $M$ is not faithfully flat.

A useful sufficient condition follows with the actual
residue maps retained. Suppose every $M_i$ is faithfully
flat and, for each maximal $\mathfrak m\subset R$, there
is a tail $i\geq i_0$ on which all maps
$M_i/\mathfrak m_iM_i\to M_j/\mathfrak m_jM_j$ are injective.
Faithful flatness at $i_0$ makes its quotient nonzero because
$R_{i_0}/\mathfrak m_{i_0}\ne0$. A nonzero element there
remains nonzero at every later stage, which is precisely
the necessary and sufficient condition just proved.

One condition on the original transitions that guarantees
this is purity of the scalar-extended maps
$$
g_{ij}:R_j\otimes_{R_i}M_i\longrightarrow M_j,
\qquad r_j\otimes u_i\longmapsto r_jf_{ij}(u_i).
$$
Here purity means that tensoring this map with every
$R_j$-module gives an injection. The formula is balanced
by semilinearity, so $g_{ij}$ is an $R_j$-linear map.
The contraction equality gives an injection of
$R_i$-modules $R_i/\mathfrak m_i\to R_j/\mathfrak m_j$.
Flatness of $M_i$ and then purity give injections
$$
(R_i/\mathfrak m_i)\otimes_{R_i}M_i
\longrightarrow (R_j/\mathfrak m_j)\otimes_{R_i}M_i
\longrightarrow (R_j/\mathfrak m_j)\otimes_{R_j}M_j.
$$
For the second map the associativity comparison sends
$(s_j+\mathfrak m_j)\otimes(r_j\otimes u_i)$ to
$(s_jr_j+\mathfrak m_j)\otimes u_i$ and has inverse
$(s_j+\mathfrak m_j)\otimes u_i\mapsto
(s_j+\mathfrak m_j)\otimes(1\otimes u_i)$; balancing and
these formulas verify both composites. Under the quotient
identifications, the composite injection is exactly
$u_i+\mathfrak m_iM_i\mapsto f_{ij}(u_i)+\mathfrak m_jM_j$.
Thus faithful flatness of the stages together with these
pure transition maps suffices.

If $R=0$, its unital action forces $M=0$. The only unital
module over that ring is zero, so its tensor functor
preserves and reflects exactness; it is faithfully flat in
the stated sense. The maximal-ideal criterion gives the
same result because there are no maximal ideals.
\end{proof}

\begin{lemma}
\label{editorial-source-lemma-flat-base-change}
Suppose that $M$ is (faithfully) flat over $R$, and that $R \to R'$
is a ring map. Then $M \otimes_R R'$ is (faithfully) flat over $R'$.
\end{lemma}

\begin{proof}
For every original $R'$-module $N$, define
$$
\begin{aligned}
\theta_N:N\otimes_{R'}(R'\otimes_R M)&\longrightarrow N\otimes_R M,
&n\otimes(r'\otimes m)&\longmapsto r'n\otimes m,\\
\eta_N:N\otimes_R M&\longrightarrow N\otimes_{R'}(R'\otimes_R M),
&n\otimes m&\longmapsto n\otimes(1\otimes m).
\end{aligned}
$$
The first formula respects the outer $R'$-balancing and the
inner $R$-balancing by the original $R'$-action on $N$.
The inverse respects $R$-balancing because
$(rn)\otimes(1\otimes m)=n\otimes(\varphi(r)\otimes m)
=n\otimes(1\otimes rm)$. Both composites fix the tensor
generators, using $r'n\otimes(1\otimes m)=n\otimes(r'\otimes m)$.
The formulas also commute with every $R'$-linear map of $N$.
The flip $M\otimes_R R'\to R'\otimes_R M$, $m\otimes r'\mapsto r'\otimes m$,
identifies the module in the statement with this one and is
its own inverse after reversing the factors. Exactness of a
complex of $R'$-modules is equivalent to exactness of its
underlying $R$-module complex, since kernels and images are
the same subsets. Apply flatness, or preservation and
reflection of exactness, for the original $M$ over $R$ and
then use the displayed natural maps. This proves both claims.
\end{proof}

\begin{lemma}
\label{editorial-source-lemma-flatness-descends}
Let $R \to R'$ be a faithfully flat ring map.
Let $M$ be a module over $R$, and set $M' = R' \otimes_R M$.
Then $M$ is flat over $R$ if and only if $M'$ is flat over $R'$.
\end{lemma}

\begin{proof}
By Lemma \ref{editorial-source-lemma-flat-base-change} we see that if $M$ is flat
then $M'$ is flat. For the converse, suppose that $M'$ is flat.
Let $N_1 \to N_2 \to N_3$ be an exact sequence of $R$-modules.
We want to show that $N_1 \otimes_R M \to N_2 \otimes_R M \to N_3 \otimes_R M$
is exact. We know that
$N_1 \otimes_R R' \to N_2 \otimes_R R' \to N_3 \otimes_R R'$ is
exact, because $R \to R'$ is flat. Flatness of $M'$ implies that
$N_1 \otimes_R R' \otimes_{R'} M'
\to N_2 \otimes_R R' \otimes_{R'} M'
\to N_3 \otimes_R R' \otimes_{R'} M'$ is exact.
The exact natural comparison for every original $N$ is
\[
\begin{split}
(N\otimes_R R')\otimes_{R'}(R'\otimes_R M)
&\longrightarrow (N\otimes_R M)\otimes_R R',\\
(n\otimes a)\otimes(b\otimes m)&\longmapsto(n\otimes m)\otimes ab.
\end{split}
\]
Its inverse sends $(n\otimes m)\otimes c$ to
$(n\otimes c)\otimes(1\otimes m)$. All $R$-balancing
relations follow from the original coefficient map $R\to R'$;
the $R'$-balancing relation moves a scalar between $a$ and
$b$ and preserves their product $ab$. These observations
make both maps well-defined. One composite fixes
$(n\otimes m)\otimes c$. The other gives
$(n\otimes ab)\otimes(1\otimes m)
=(n\otimes a)\otimes(b\otimes m)$ by $R'$-balancing.
Their formulas commute with every map of $N$. Thus the
exact complex above is the tensor with $R'$ of the original
complex $N_1\otimes_R M\to N_2\otimes_R M\to N_3\otimes_R M$.
Faithful flatness of $R'$ over $R$ reflects exactness, so
this original complex is exact, as required.
\end{proof}

\begin{lemma}
\label{editorial-source-lemma-flatness-descends-more-general}
Let $R$ be a ring. Let $S \to S'$ be a flat map of $R$-algebras.
Let $M$ be a module over $S$, and set $M' = S' \otimes_S M$.
\begin{enumerate}
\item If $M$ is flat over $R$, then $M'$ is flat over $R$.
\item If $S \to S'$ is faithfully flat, then $M$ is flat
over $R$ if and only if $M'$ is flat over $R$.
\end{enumerate}
\end{lemma}

\begin{proof}
Let $f:N\to N'$ be an injection of $R$-modules and put
$K=\Ker(f\otimes\mathrm{id}_M)$. The map
$f\otimes\mathrm{id}_M:N\otimes_R M\to N'\otimes_R M$
is $S$-linear, with the $S$-action on its $M$ factor, so
$K$ is an $S$-module. Tensoring its defining exact sequence
with the flat $S$-module $S'$ identifies $K\otimes_S S'$
with the kernel after that tensor operation. For every $N$
the precise comparison is
\[
\begin{split}
(N\otimes_R M)\otimes_S S'&\longrightarrow
N\otimes_R(S'\otimes_S M),\\
(n\otimes m)\otimes s'&\longmapsto n\otimes(s'\otimes m),
\end{split}
\]
with inverse $n\otimes(s'\otimes m)\mapsto(n\otimes m)\otimes s'$.
The $S$-balancing relations agree on the two sides by moving
the original $S$-scalar between $m$ and $s'$; the $R$-balancing
relations agree through $R\to S\to S'$. The composites fix
the generators, and the formulas commute with $f$. Consequently
the kernel identification is exactly
$$
K\otimes_S S'\ \cong\
\Ker(N\otimes_R M'\longrightarrow N'\otimes_R M').
$$
If $M$ is flat over $R$, then $K=0$, so the right kernel
vanishes for every injection $f$ and Lemma \ref{editorial-source-lemma-flat}
gives flatness of $M'$ over $R$. Conversely, if $M'$ is
flat, then $K\otimes_S S'=0$. When $S'$ is faithfully flat
over $S$, apply reflection of exactness to the complex
$0\to K\to0$: its tensor complex is exact, hence $K=0$.
The same criterion now gives flatness of $M$ over $R$.
\end{proof}

\begin{lemma}
\label{editorial-source-lemma-flat-permanence}
Let $R \to S$ be a ring map. Let $M$ be an $S$-module.
If $M$ is flat as an $R$-module and faithfully flat as an $S$-module,
then $R\to S$ is flat.
More generally, whenever $M$ is faithfully flat over $S$,
the map $R\to S$ is flat (respectively faithfully flat)
if and only if $M$ is flat (respectively faithfully flat)
as an $R$-module.
\end{lemma}

\begin{proof}
For every original $R$-module $N$, the comparison
$$
(N\otimes_R S)\otimes_S M\longrightarrow N\otimes_R M,
\qquad (n\otimes s)\otimes m\longmapsto n\otimes sm
$$
has inverse $n\otimes m\mapsto(n\otimes1)\otimes m$.
Balancing over $S$ preserves the product $sm$, and balancing
over $R$ is compatible with the restricted action on $M$.
Both composites fix the generators; for one of them use
$(n\otimes1)\otimes sm=(n\otimes s)\otimes m$.
These maps are natural in $N$.

If $M$ is flat over $R$, an exact complex
$N_1\to N_2\to N_3$ stays exact after tensoring with $M$.
The comparison identifies this with the result of first
tensoring with $S$ and then with $M$ over $S$. Faithful
flatness over $S$ reflects exactness, so the first tensor
complex is exact. This proves the original conclusion that
$R\to S$ is flat. The converse under the same faithful
$S$-module hypothesis follows from
Lemma \ref{editorial-source-lemma-composition-flat}.

If $M$ is faithfully flat over $R$, take any original complex
$N_1\to N_2\to N_3$ whose tensor with $S$ is exact. Tensoring
that complex with the flat $S$-module $M$ and using the
comparison shows its tensor with $M$ over $R$ is exact.
Faithful flatness over $R$ reflects this to exactness of the
original complex. Together with flatness just proved, this
makes $R\to S$ faithfully flat. Conversely, if $R\to S$
is faithfully flat, composition with the faithfully flat
$S$-module $M$ gives faithful flatness over $R$, again by
Lemma \ref{editorial-source-lemma-composition-flat}.
\end{proof}

\noindent
Let $R$ be a ring.
Let $M$ be an $R$-module.
Let $\sum_{i=1}^n f_ix_i=0$ be a finite relation in $M$,
with $f_i\in R$, $x_i\in M$, and $n\geq0$.
We say the relation $\sum f_i x_i$
is {\it trivial} if there exist an integer $m \geq 0$,
elements $y_j \in M$, $j = 1, \ldots, m$, and elements $a_{ij} \in R$,
$i = 1, \ldots, n$, $j = 1, \ldots, m$ such that
$$
x_i = \sum\nolimits_j a_{ij} y_j, \forall i,
\quad\text{and}\quad
0 = \sum\nolimits_i f_ia_{ij}, \forall j.
$$

\begin{lemma}[Equational criterion of flatness]
\label{editorial-source-lemma-flat-eq}
A module $M$ over $R$ is flat if and only if
every relation in $M$ is trivial.
\end{lemma}

\begin{proof}
Assume $M$ is flat and let $\sum f_i x_i = 0$ be a relation in $M$.
Let $I = (f_1, \ldots, f_n)$, and let
$K = \Ker(R^n \to I, (a_1, \ldots, a_n) \mapsto \sum_i a_i f_i)$.
So we have the short exact sequence
$0 \to K \to R^n \to I \to 0$. Then $\sum f_i \otimes x_i$
is an element of $I \otimes_R M$ which maps
to zero in $R \otimes_R M = M$. By flatness
$\sum f_i \otimes x_i$ is zero in $I \otimes_R M$.
Thus there exists an element of $K \otimes_R M$ mapping
to $\sum e_i \otimes x_i \in R^n \otimes_R M$ where $e_i$
is the $i$th basis element of $R^n$.
Write this element as a finite sum
$\sum_{j=1}^m k_j\otimes y_j$, and write each original
$k_j\in K\subseteq R^n$ as $\sum_{i=1}^n a_{ij}e_i$.
Membership in $K$ gives $\sum_{i=1}^n f_i a_{ij}=0$
for every $j$. Under the coordinate isomorphism
$R^n\otimes_R M\to M^n$, $e_i\otimes y$ maps to the
vector having $y$ in coordinate $i$ and zero elsewhere;
its inverse sends $(x_1,\ldots,x_n)$ to
$\sum_i e_i\otimes x_i$. Comparing all coordinates of
the original lifted tensor gives
$x_i=\sum_{j=1}^m a_{ij}y_j$ for every $i$, exactly the
two defining conditions. The sum may be empty; when $n=0$
there are no coordinates and the empty relation is handled
with $m=0$.

\medskip\noindent
Assume every relation is trivial, let $I$
be a finitely generated ideal, and let $x = \sum f_i \otimes x_i$
be an element of $I \otimes_R M$ mapping to zero in $R \otimes_R M = M$.
This just means exactly that $\sum f_i x_i$ is a relation in
$M$. And the fact that it is trivial implies easily that
$x$ is zero, because
$$
x
=
\sum f_i \otimes x_i
=
\sum f_i \otimes \left(\sum a_{ij}y_j\right)
=
\sum \left(\sum f_i a_{ij}\right) \otimes y_j
=
0
$$
The sums are finite, so the middle equality follows by
additivity and the original balancing relation
$f_i\otimes a_{ij}y_j=(f_i a_{ij})\otimes y_j$ in
$I\otimes_R M$. Thus the map $I\otimes_R M\to M$ is
injective for every finitely generated ideal $I$, and
Lemma \ref{editorial-source-lemma-flat} proves that $M$ is flat.
\end{proof}

\begin{lemma}
\label{editorial-source-lemma-flat-finite-simultaneous-relations}
An $R$-module $M$ is flat if and only if it has the following
property. For every matrix $(f_{\ell i})\in\operatorname{Mat}_{r\times n}(R)$
and elements $x_1,\ldots,x_n\in M$ satisfying
$\sum_{i=1}^n f_{\ell i}x_i=0$ for every $1\leq\ell\leq r$,
there are an integer $m\geq0$, elements $y_1,\ldots,y_m\in M$
and coefficients $a_{ij}\in R$ such that
$$
x_i=\sum_{j=1}^m a_{ij}y_j\quad(1\leq i\leq n),\qquad
\sum_{i=1}^n f_{\ell i}a_{ij}=0
\quad(1\leq\ell\leq r,\ 1\leq j\leq m).
$$
Here $r,n$ may be zero; all matrices and sums are finite.
\end{lemma}

\begin{proof}
Keep the original matrix and let
$\psi:R^n\to R^r$ send the $i$th basis vector to the
column $(f_{1i},\ldots,f_{ri})$. Put $K=\Ker\psi$.
Flatness makes
$K\otimes_R M\to R^n\otimes_R M\to R^r\otimes_R M$
exact. The coordinate tensor isomorphisms in the preceding
proof identify its last map with the original matrix
acting on $M^n$. Thus the given vector is the image of
a finite sum $\sum_{j=1}^m k_j\otimes y_j$.
Write $k_j=\sum_i a_{ij}e_i$. All $r$ original equations
for $k_j\in K$ are precisely the second displayed family,
and coordinate equality gives the first family. No finite
generation assumption on $K$ is used. Conversely, the case
$r=1$ gives the equational criterion
of Lemma \ref{editorial-source-lemma-flat-eq}, hence flatness. When $r=0$
one may take $m=n$, $a_{ij}=\delta_{ij}$ and $y_j=x_j$;
when $n=0$ take $m=0$. These choices verify the boundary
cases with the original empty matrices retained.
\end{proof}

\begin{lemma}
\label{editorial-source-lemma-flat-tor-zero}
Suppose that $R$ is a ring, that $0\to M''\to M'\to M\to0$
is a short exact sequence, and that $N$ is an $R$-module. If $M$ is flat
then $N \otimes_R M'' \to N \otimes_R M'$ is injective, i.e., the
sequence
$$
0 \to N \otimes_R M'' \to N \otimes_R M' \to N \otimes_R M \to 0
$$
is a short exact sequence.
\end{lemma}

\begin{proof}
Retain a surjection $q:F=R^{(I)}\to N$ with kernel
$K$ and inclusion $j:K\to F$. Write $u:M''\to M'$ and
$v:M'\to M$ for the original maps. The coordinate
isomorphism $L\otimes_R F\to L^{(I)}$ for any $L$ sends
$\ell\otimes\sum_i r_ie_i$ to $(r_i\ell)_i$ and has
inverse $(\ell_i)_i\mapsto\sum_i\ell_i\otimes e_i$;
all supports are finite. The rows and columns of the
following diagram are induced by $u,v,j,q$:
$$
\begin{matrix}
 & & 0 & & 0 & & 0 & & \\
 & & \uparrow & & \uparrow & & \uparrow & & \\
 & & M''\otimes_R N & \to & M' \otimes_R N & \to & M \otimes_R N & \to & 0 \\
 & & \uparrow & & \uparrow & & \uparrow & & \\
0 & \to & (M'')^{(I)} & \to & (M')^{(I)} & \to & M^{(I)} & \to & 0 \\
 & & \uparrow & & \uparrow & & \uparrow & & \\
 & & M''\otimes_R K & \to & M' \otimes_R K & \to & M \otimes_R K & \to & 0 \\
 & & & & & & \uparrow & & \\
 & & & & & & 0 & &
\end{matrix}
$$
Each column ending in zero is right exact by tensoring
$K\to F\to N\to0$. The bottom and top rows are right
exact by tensoring $M''\to M'\to M\to0$. The middle row
is short exact coordinate by coordinate, including
surjectivity since a finitely supported tuple has finitely
many lifts. The bottom right vertical map is injective
because $M$ is flat and $K\to F$ is injective.
Every square commutes on its tensor generators.

Here is the full chase proving the remaining injection.
Let $x\in M''\otimes_R N$ map to zero in $M'\otimes_R N$.
Lift it to $a\in(M'')^{(I)}$. The image $u^{(I)}(a)$
has zero image in $M'\otimes_R N$, so it is the image
of some $b\in M'\otimes_R K$. The image of $b$ in
$M\otimes_R K$ maps to zero in $M^{(I)}$, because
$v^{(I)}u^{(I)}(a)=0$. Bottom right injectivity forces
that image to be zero. Bottom-row exactness supplies
$c\in M''\otimes_R K$ mapping to $b$.
Its image $d\in(M'')^{(I)}$ satisfies
$u^{(I)}(d)=u^{(I)}(a)$ by commutativity. Middle-row
injectivity gives $d=a$. As $d$ comes from
$M''\otimes_R K$, its image in $M''\otimes_R N$ is zero,
so $x=0$. This proves the asserted injection.
Finally the flips $L\otimes_R N\to N\otimes_R L$,
$\ell\otimes n\mapsto n\otimes\ell$, are mutual inverses
and commute with $u,v$, so the diagram conclusion is
exactly the sequence in the statement.
\end{proof}

\begin{lemma}
\label{editorial-source-lemma-flat-quotient-cyclic-tensor-criterion}
In the original short exact sequence
$0\to M''\xrightarrow{u}M'\xrightarrow{v}M\to0$,
assume $M'$ is flat. The following are equivalent:
\begin{enumerate}
\item $M$ is flat.
\item For every $R$-module $N$, the map
$N\otimes_R M''\to N\otimes_R M'$ is injective.
\item For every finitely generated ideal $I\subseteq R$,
the map $M''/IM''\to M'/IM'$ induced by $u$ is injective.
\end{enumerate}
\end{lemma}

\begin{proof}
The first implication is Lemma \ref{editorial-source-lemma-flat-tor-zero}.
The second follows by taking $N=R/I$ and using the
quotient-tensor maps $(r+I)\otimes m\mapsto rm+IM$
and $m+IM\mapsto(1+I)\otimes m$, on each original module.

Assume the third condition and retain any finite original
relation $\sum_{i=1}^n f_ix_i=0$ in $M$.
Choose lifts $\widetilde x_i\in M'$ with
$v(\widetilde x_i)=x_i$, and set $I=(f_1,\ldots,f_n)$.
The element $\sum_i f_i\widetilde x_i$ has zero image
under $v$, so it equals $u(w)$ for a unique $w\in M''$.
Its class in $M'/IM'$ is zero. Injectivity of the specified
quotient map gives $w\in IM''$.
Write $w=\sum_{\ell=1}^q c_\ell z_\ell$ with
$c_\ell\in I$, $z_\ell\in M''$, and express each original
$c_\ell=\sum_{i=1}^n f_i d_{i\ell}$.
Putting $w_i=\sum_{\ell=1}^q d_{i\ell}z_\ell$ gives the
exact equality $w=\sum_i f_iw_i$.
The corrected lifts
$\widehat x_i=\widetilde x_i-u(w_i)$ therefore satisfy
$$
\sum_i f_i\widehat x_i
=u(w)-u\left(\sum_i f_iw_i\right)=0,
\qquad v(\widehat x_i)=x_i.
$$
By flatness of the original $M'$ and
Lemma \ref{editorial-source-lemma-flat-eq}, there are $y_j\in M'$ and
$a_{ij}\in R$ with
$\widehat x_i=\sum_j a_{ij}y_j$ and
$\sum_i f_i a_{ij}=0$ for every $j$.
Applying $v$ gives $x_i=\sum_j a_{ij}v(y_j)$ with the
same coefficients and the same vanishing equations.
Thus every original relation in $M$ is trivial and
$M$ is flat by that criterion. The empty relation
requires no lifts and is already trivial.
\end{proof}


\begin{lemma}
\label{editorial-source-lemma-flat-ses}
Suppose that $0 \to M' \to M \to M'' \to 0$ is
a short exact sequence of $R$-modules.
If $M'$ and $M''$ are flat so is $M$.
If $M$ and $M''$ are flat so is $M'$.
\end{lemma}

\begin{proof}
We will use the criterion that a module $N$ is flat if for
every ideal $I \subset R$ the map $N \otimes_R I \to N$ is injective,
see Lemma \ref{editorial-source-lemma-flat}.
Consider an ideal $I \subset R$.
Consider the diagram
$$
\begin{matrix}
0 & \to & M' & \to & M & \to & M'' & \to & 0 \\
& & \uparrow & & \uparrow & & \uparrow & & \\
& & M'\otimes_R I & \to & M \otimes_R I & \to & M''\otimes_R I & \to & 0
\end{matrix}
$$
The vertical maps are $m\otimes a\mapsto am$ for the
original ideal $I$, so both squares commute on generators.
The upper row is the given short exact sequence and the
lower row is right exact. For the first assertion suppose
$M'$ and $M''$ are flat, and let $x\in M\otimes_R I$
map to zero in $M$. Its image in $M''\otimes_R I$ maps
to zero in $M''$ and hence is zero by flatness of $M''$.
Lower-row exactness lifts $x$ to $y\in M'\otimes_R I$.
The image of $y$ in $M'$ maps to zero in $M$, so it is
zero by upper-row injectivity. Flatness of $M'$ now gives
$y=0$, hence $x=0$. Lemma \ref{editorial-source-lemma-flat} proves that
$M$ is flat.

For the second assertion suppose $M$ and $M''$ are flat.
Lemma \ref{editorial-source-lemma-flat-tor-zero}, applied with quotient
$M''$ and tensor module $I$, makes the lower left map
injective. If $y\in M'\otimes_R I$ maps to zero in $M'$,
its image $x\in M\otimes_R I$ maps to zero in $M$.
Flatness of $M$ gives $x=0$, and lower left injectivity
gives $y=0$. The same ideal criterion proves flatness
of the original $M'$.
\end{proof}

\begin{lemma}
\label{editorial-source-lemma-easy-ff}
Let $R$ be a ring.
Let $M$ be an $R$-module.
The following are equivalent:
\begin{enumerate}
\item $M$ is faithfully flat, and
\item $M$ is flat and for all $R$-module homomorphisms $\alpha : N \to N'$
we have $\alpha = 0$ if and only if $\alpha \otimes \text{id}_M = 0$.
\end{enumerate}
\end{lemma}

\begin{proof}
Faithful flatness includes flatness. It also implies that
$L\otimes_R M=0$ forces $L=0$, by reflecting exactness
of $0\to L\to0$. For an original map $\alpha:N\to N'$,
retain $K=\Ker\alpha$. Flatness makes
$0\to K\otimes_R M\to N\otimes_R M\to N'\otimes_R M$
exact. If $\alpha\otimes\mathrm{id}_M=0$, the first
injection is also surjective. Right exactness of tensoring
$K\to N\to N/K\to0$ then gives $(N/K)\otimes_R M=0$.
The vanishing reflection just proved gives $N/K=0$, so
$K=N$ and $\alpha=0$. Conversely a zero map tensors
to a zero map by its formula on generators. This proves
the first implication with the original kernel retained.

Assume the second condition. Flatness already preserves
exactness. To prove reflection, take an original complex
$N_1\xrightarrow{d_1}N_2\xrightarrow{d_2}N_3$ whose
tensor complex is exact, and put $J=\operatorname{Im}d_1$.
For each $x\in\Ker d_2$, keep the original map
$\alpha:R\to N_2/J$, $r\mapsto rx+J$.
For every $m\in M$, the element $x\otimes m$ belongs to
$\Ker(d_2\otimes\mathrm{id}_M)$, hence to the image of
$d_1\otimes\mathrm{id}_M$. The quotient map $N_2\to N_2/J$
tensored with $M$ kills that image, so
$(x+J)\otimes m=0$. Therefore on every tensor generator
$$
(\alpha\otimes\mathrm{id}_M)(r\otimes m)
=(rx+J)\otimes m=(x+J)\otimes rm=0.
$$
The assumed reflection of zero morphisms gives $\alpha=0$,
and evaluation at $1$ gives $x\in J$. Thus
$\Ker d_2\subseteq\operatorname{Im}d_1$; the reverse
inclusion holds because the original sequence is a complex.
It is exact, proving faithful flatness.
\end{proof}

\begin{lemma}
\label{editorial-source-lemma-ff}
\begin{slogan}
A flat module is faithfully flat if and only if it has nonzero fibers.
\end{slogan}
Let $M$ be a flat $R$-module.
The following are equivalent:
\begin{enumerate}
\item $M$ is faithfully flat,
\item for every nonzero $R$-module $N$, the tensor product $M \otimes_R N$
is nonzero,
\item for all $\mathfrak p \in \Spec(R)$
the tensor product $M \otimes_R \kappa(\mathfrak p)$ is nonzero, and
\item for all maximal ideals $\mathfrak m$ of $R$
the tensor product $M \otimes_R \kappa(\mathfrak m) = M/{\mathfrak m}M$
is nonzero.
\end{enumerate}
\end{lemma}

\begin{proof}
Assume $M$ faithfully flat and $N \not = 0$. By Lemma \ref{editorial-source-lemma-easy-ff}
the nonzero map $1 : N \to N$ induces a nonzero map
$M \otimes_R N \to M \otimes_R N$, so $M \otimes_R N \not = 0$.
Thus (1) implies (2). The implications (2) $\Rightarrow$ (3) $\Rightarrow$ (4)
are immediate.

\medskip\noindent
Assume (4), and retain a complex
$N_1\xrightarrow{d_1}N_2\xrightarrow{d_2}N_3$ whose
tensor complex is exact. Write
$K=\Ker d_2$, $J=\operatorname{Im}d_1\subseteq K$,
and keep its original cohomology $H=K/J$.
Flatness embeds $J\otimes_R M$ and $K\otimes_R M$ in
$N_2\otimes_R M$. Tensoring the surjection $N_1\to J$
identifies the first image with
$\operatorname{Im}(d_1\otimes\mathrm{id}_M)$; tensoring
$0\to K\to N_2\to N_3$ identifies the second image with
$\operatorname{Ker}(d_2\otimes\mathrm{id}_M)$.
The tensor complex is exact, so these submodules coincide.
Right exactness applied to $J\to K\to H\to0$ gives
$H\otimes_R M=0$ with these exact original maps.

Take any $x\in H$ and set $I=\{f\in R:fx=0\}$.
The map $R/I\to H$, $f+I\mapsto fx$, is well-defined
and injective by the definition of $I$.
Flatness makes $(R/I)\otimes_R M\to H\otimes_R M=0$
injective. The quotient-tensor isomorphism
$(f+I)\otimes m\mapsto fm+IM$ with inverse
$m+IM\mapsto(1+I)\otimes m$ therefore gives $M/IM=0$.
If $I$ were proper, choose a maximal ideal
$\mathfrak m\supseteq I$. The quotient map
$M/IM\to M/\mathfrak mM$, $m+IM\mapsto m+\mathfrak mM$,
is well-defined and surjective. Its source is zero, forcing
$M/\mathfrak mM=0$, contrary to (4). Thus $I=R$,
$x=1x=0$, and $H=0$. This proves exactness reflection.
For the zero ring all modules are zero and all prime and
maximal-ideal conditions are vacuous, so the same equivalence
also covers that case.
\end{proof}

\begin{lemma}
\label{editorial-source-lemma-ff-rings}
Let $R \to S$ be a flat ring map.
The following are equivalent:
\begin{enumerate}
\item $R \to S$ is faithfully flat,
\item the induced map on $\Spec$ is surjective, and
\item any closed point $x \in \Spec(R)$ is
in the image of the map $\Spec(S) \to \Spec(R)$.
\end{enumerate}
\end{lemma}

\begin{proof}
For the original prime $\mathfrak p\subset R$, the fibre
ring $S\otimes_R\kappa(\mathfrak p)$ is nonzero exactly
when there is a prime of $S$ contracting to $\mathfrak p$,
by Remark \ref{algebra-remark-fundamental-diagram} and
Lemma \ref{editorial-source-lemma-in-image}. The prime correspondence is
through the quotient by $\mathfrak p$ followed by inversion
of the original $R\setminus\mathfrak p$; its primes
contract to primes of $S$ with precisely that contraction.
Thus condition (3) of Lemma \ref{editorial-source-lemma-ff} says exactly
that every prime is in the spectrum image. Its condition
(4) says exactly that every maximal ideal is in the image,
and maximal ideals are the closed points of $\Spec(R)$.
Since $S$ is flat over $R$, that lemma proves all three
equivalences asserted here.
\end{proof}

\begin{lemma}
\label{editorial-source-lemma-local-flat-ff}
A flat local ring homomorphism of local rings is faithfully flat.
\end{lemma}

\begin{proof}
Write $\varphi:(R,\mathfrak m)\to(S,\mathfrak n)$ for
the original local homomorphism. Its defining local
condition is $\varphi^{-1}(\mathfrak n)=\mathfrak m$.
Thus the unique closed point of $\Spec(R)$ is the image
of the prime $\mathfrak n$ of $S$. Flatness and
Lemma \ref{editorial-source-lemma-ff-rings} give faithful flatness.
\end{proof}

\noindent
Flatness meshes well with localization.

\begin{lemma}
\label{editorial-source-lemma-flat-localization}
Let $R$ be a ring. Let $S \subset R$ be a multiplicative subset.
\begin{enumerate}
\item The localization $S^{-1}R$ is a flat $R$-algebra.
\item If $M$ is an $S^{-1}R$-module, then $M$ is a flat $R$-module
if and only if $M$ is a flat $S^{-1}R$-module.
\item Suppose $M$ is an $R$-module. Then
$M$ is a flat $R$-module if and only if $M_{\mathfrak p}$ is a flat
$R_{\mathfrak p}$-module for all primes $\mathfrak p$ of $R$.
\item Suppose $M$ is an $R$-module. Then $M$ is a flat $R$-module if
and only if $M_{\mathfrak m}$ is a flat
$R_{\mathfrak m}$-module for all maximal ideals $\mathfrak m$ of $R$.
\item Suppose $R \to A$ is a ring map, $M$ is an $A$-module,
and $g_1, \ldots, g_m \in A$ are elements generating the unit
ideal of $A$. Then $M$ is flat over $R$ if and only if each localization
$M_{g_i}$ is flat over $R$.
\item Suppose $R \to A$ is a ring map, and $M$ is an $A$-module.
Then $M$ is a flat $R$-module if and only if the localization
$M_{\mathfrak q}$ is a flat $R_{\mathfrak p}$-module
(with $\mathfrak p$ the prime of $R$ lying under $\mathfrak q$)
for all primes $\mathfrak q$ of $A$.
\item Suppose $R \to A$ is a ring map, and $M$ is an $A$-module.
Then $M$ is a flat $R$-module if and only if the localization
$M_{\mathfrak m}$ is a flat $R_{\mathfrak p}$-module
(with $\mathfrak p$ the inverse image of $\mathfrak m$ under $R\to A$)
for all maximal ideals $\mathfrak m$ of $A$.
\end{enumerate}
\end{lemma}

\begin{proof}
Retain every original coefficient map; in (5)--(7) write it
as $\varphi:R\to A$. All contractions below are inverse
images under this map, which need not be injective.

We first recall the exact fraction facts used in the proof.
For a ring $B$, a multiplicative set $W\subseteq B$ and a
$B$-module $L$, the equality $\ell/w=\ell'/w'$ means
that $v(w'\ell-w\ell')=0$ for some original $v\in W$.
Thus $\ell/w=0$ means that some $v\in W$ kills $\ell$.
If $L'\to L$ is injective, this equality for an element
of $L'$ holds in $L'$ whenever it holds in $L$, so
localization preserves injections. For a map $u:L\to N$,
an element $\ell/w$ in the localized kernel has
$v u(\ell)=0$ for some $v\in W$. Then $v\ell\in\Ker u$
and $\ell/w=(v\ell)/(vw)$, proving that localization
preserves kernels. Its image consists exactly of fractions
$u(\ell)/w$, and a surjection remains surjective by lifting
the numerator. These facts prove exactness with the
original fraction witnesses retained.

For later zero detection, if $0\ne\ell\in L$ its
annihilator is a proper ideal of $B$ and is contained in
a maximal ideal $\mathfrak m$. Then $\ell/1\ne0$ in
$L_{\mathfrak m}$, since an element outside $\mathfrak m$
killing $\ell$ would belong to that annihilator. Hence
$L=0$ if all its maximal localizations vanish, as in
Lemma \ref{algebra-lemma-characterize-zero-local}.

For an original $R$-module $N$, the $A$-action on
$N\otimes_R M$ is $a(n\otimes m)=n\otimes am$.
For every multiplicative $W\subseteq A$, there are
mutually inverse maps
\[
\begin{split}
W^{-1}(N\otimes_R M)&\longrightarrow N\otimes_R W^{-1}M,\\
\frac{\sum_{\ell=1}^d n_\ell\otimes m_\ell}{w}
&\longmapsto\sum_{\ell=1}^d n_\ell\otimes\frac{m_\ell}{w},
\end{split}
\]
and $n\otimes(m/w)\mapsto(n\otimes m)/w$.
For the forward map start with $n\otimes m\mapsto
n\otimes(m/1)$, which is $A$-linear and $R$-balanced;
every element of $W$ acts invertibly on its target, so
it extends to the localization with the stated formula.
For the inverse, if $v(w'm-wm')=0$, then
$v(w'(n\otimes m)-w(n\otimes m'))=0$ by the original
$A$-action. This proves independence of the fraction.
The $R$-balancing relation is
$rn\otimes m=n\otimes\varphi(r)m$, and is preserved by
the formula. Additivity follows by common denominators.
Both composites fix the displayed generators and their
finite sums. The maps are natural in $N$ and $M$, since
they apply their homomorphisms only to the numerators.
Consequently, if $M$ is flat over $R$, then $W^{-1}M$
is flat over $R$: localize the injection obtained by
tensoring any injection of $R$-modules with $M$ and use
these natural isomorphisms.

For a prime $\mathfrak q\subset A$, put
$\mathfrak p=\varphi^{-1}(\mathfrak q)$. The map and action are
$$
R_{\mathfrak p}\longrightarrow A_{\mathfrak q},\quad
\frac rs\longmapsto\frac{\varphi(r)}{\varphi(s)},
\qquad
\frac rs\cdot\frac ma=\frac{\varphi(r)m}{\varphi(s)a}.
$$
All denominators are valid because $s\notin\mathfrak p$
is equivalent to $\varphi(s)\notin\mathfrak q$.
The exact two-sided comparison for every original $N$ is
$$
(N\otimes_R M)_{\mathfrak q}
\ \cong\ N_{\mathfrak p}\otimes_{R_{\mathfrak p}}M_{\mathfrak q},
\qquad
\frac{n\otimes m}{a}\longmapsto\frac n1\otimes\frac ma,
$$
with inverse
$$
\frac ns\otimes\frac ma\longmapsto
\frac{n\otimes m}{\varphi(s)a}.
$$
If $n/s=n'/s'$, choose $u\notin\mathfrak p$ with
$u(s'n-sn')=0$. The two proposed inverse images have
common denominator $\varphi(s)\varphi(s')a$ and numerator
$(s'n-sn')\otimes m$. Multiplication by
$\varphi(u)\notin\mathfrak q$ kills it, by $R$-balancing.
If $m/a=m'/a'$, choose $b\notin\mathfrak q$ with
$b(a'm-am')=0$. Using common denominator $\varphi(s)aa'$
gives numerator $n\otimes(a'm-am')$, killed by $b$.
For $r/t\in R_{\mathfrak p}$ the two balancing images are
$$
\frac{rn\otimes m}{\varphi(t)\varphi(s)a}
=\frac{n\otimes\varphi(r)m}{\varphi(s)\varphi(t)a},
$$
so the inverse is well-defined and additive.
The forward map is the localization of
$n\otimes m\mapsto(n/1)\otimes(m/1)$; elements outside
$\mathfrak q$ act invertibly on its target.
One composite fixes $(n\otimes m)/a$, while the other
sends $(n/s)\otimes(m/a)$ to
$(n/1)\otimes(m/(\varphi(s)a))$, which equals the
original tensor by balancing with $1/s$.
The formulas prove naturality for the original maps.

We also retain the exact ideal descent used in the source.
For any ideal $I\subseteq R_{\mathfrak p}$, put
$J=\{r\in R:r/1\in I\}$. Then $J_{\mathfrak p}=I$:
if $j\in J$, its fraction $j/s=(1/s)(j/1)$ is in $I$;
if $a/s\in I$, then $a/1=(s/1)(a/s)\in I$, so $a\in J$.
The localized map $J_{\mathfrak p}\to R_{\mathfrak p}$ is
injective by localization of $J\hookrightarrow R$.
For specified finite generators $a_\ell/s_\ell$ of $I$,
the smaller specified ideal $J_0=(a_1,\ldots,a_n)$ also
has $(J_0)_{\mathfrak p}=I$, because each $s_\ell/1$ is
a unit and both identities above apply to that generator.
This does not assert finite generation of the contraction $J$.

We now prove (1). Keep the original multiplicative set
$S\subseteq R$ and put $B=S^{-1}R$.
For every $N$ the maps
$$
N\otimes_R B\longrightarrow S^{-1}N,\quad
n\otimes(r/s)\longmapsto rn/s,
\qquad
n/s\longmapsto n\otimes(1/s)
$$
are inverse. To check the inverse on fraction equality,
if $u(s'n-sn')=0$ then its difference of images is
$$
(s'n-sn')\otimes\frac1{ss'}
=u(s'n-sn')\otimes\frac1{uss'}=0.
$$
The forward map respects equal fractions for the same
reason and respects the original $R$-balancing relation.
The two composites fix the generators, since
$rn\otimes(1/s)=n\otimes(r/s)$.
Localization preserves injections, so these natural maps
and Lemma \ref{editorial-source-lemma-flat} prove that $B$ is flat over $R$.

For (2), if $M$ is flat over $B$, then (1) and
Lemma \ref{editorial-source-lemma-composition-flat} imply flatness over $R$.
The receiving tensor comparison can also be written as
$$
(S^{-1}N)\otimes_B M\longrightarrow N\otimes_R M,
\qquad (n/s)\otimes m\longmapsto n\otimes(1/s)m,
$$
with inverse $n\otimes m\mapsto(n/1)\otimes m$.
For example if $u(s'n-sn')=0$, the difference between
the first images is
$(s'n-sn')\otimes(1/(ss'))m
=u(s'n-sn')\otimes(1/(uss'))m=0$.
Balancing by $r/t\in B$ agrees on both sides after
moving $r$ and the inverse action of $t$ to the $M$ factor.
Both composites are identities by these same relations.

Conversely let $M$ be a $B$-module flat over $R$.
For any $B$-module $P$, the quotient map
$P\otimes_R M\to P\otimes_B M$ has inverse
$p\otimes_B m\mapsto p\otimes_R m$.
To verify the additional balancing needed for this inverse,
for each original $s\in S$ use
$$
(1/s)p\otimes_R m
=(1/s)p\otimes_R s((1/s)m)
=p\otimes_R(1/s)m.
$$
Together with $R$-balancing this gives balancing by every
$r/s\in B$. The composites fix all generators.
An injection of $B$-modules is an injection of underlying
$R$-modules, so these natural maps and flatness over $R$
give the injection criterion over $B$. This proves (2).

Next prove (6) and (7), retaining the original ideal maps.
Suppose $M_{\mathfrak m}$ is flat over $R_{\mathfrak p}$
for each maximal $\mathfrak m\subset A$, where
$\mathfrak p=\varphi^{-1}(\mathfrak m)$.
For an ideal $I\subseteq R$, consider the original
$A$-linear map
$$
\eta:I\otimes_R M\longrightarrow M,
\qquad a\otimes m\longmapsto\varphi(a)m,
\qquad K=\Ker\eta.
$$
At $\mathfrak m$ the comparison already proved gives
$$
(I\otimes_R M)_{\mathfrak m}
\cong I_{\mathfrak p}\otimes_{R_{\mathfrak p}}M_{\mathfrak m}.
$$
It identifies $\eta_{\mathfrak m}$ with the multiplication
map $(a/s)\otimes(m/t)\mapsto\varphi(a)m/(\varphi(s)t)$,
so that map is injective by the asserted local flatness
and the inclusion $I_{\mathfrak p}\hookrightarrow R_{\mathfrak p}$.
Exactness of localization gives $K_{\mathfrak m}=0$.
Zero detection at all maximal ideals of $A$ gives $K=0$,
and Lemma \ref{editorial-source-lemma-flat} proves flatness over $R$.

Conversely suppose $M$ is flat over $R$ and retain any
prime $\mathfrak q\subset A$ with
$\mathfrak p=\varphi^{-1}(\mathfrak q)$.
For an ideal $I\subseteq R_{\mathfrak p}$ take its original
contraction $J\subseteq R$ above, so $J_{\mathfrak p}=I$.
Flatness makes $J\otimes_R M\to M$ injective. Localizing
at $\mathfrak q$ preserves its injectivity, and the exact
comparison identifies it with
$I\otimes_{R_{\mathfrak p}}M_{\mathfrak q}\to M_{\mathfrak q}$.
The ideal criterion proves flatness over $R_{\mathfrak p}$.
This proves necessity at every prime and sufficiency at
all maximal ideals. Since the former condition includes
the latter, (6) and (7) both follow. A maximal ideal of
$A$ may contract to a nonmaximal prime of $R$; no step
replaces the original contracted prime by a maximal one.

Taking $A=R$ and $\varphi=\mathrm{id}_R$ in (6) and (7)
gives (3) and (4), respectively, with the same module $M$
and the same prime and maximal localizations.

For (5) we prove the stronger arbitrary-family statement:
if $(g_\lambda)_{\lambda\in\Lambda}$ is any family in $A$
such that for every maximal $\mathfrak m$ with
$M_{\mathfrak m}\ne0$ some $g_\lambda\notin\mathfrak m$,
then $M$ is flat over $R$ if and only if every
$M_{g_\lambda}$ is flat over $R$.

First, for each original $g\in A$, multiplication by $g$
gives a directed system of the original $R$-module $M$,
indexed by $n\in\mathbf Z_{\geq0}$. Its colimit is
$$
\colim(M\xrightarrow{g}M\xrightarrow{g}\cdots)
\longrightarrow M_g,\qquad [m\text{ at }n]\longmapsto m/g^n.
$$
Compatibility is $gm/g^{n+1}=m/g^n$.
Every fraction occurs. A fraction $m/g^n$ is zero exactly
when some $g^k m=0$, which is exactly eventual vanishing
of its original representative. This proves injectivity
as well as surjectivity; the inverse sends a fraction to
that class, independently of its representative by the
same equality test. Lemma \ref{editorial-source-lemma-colimit-flat} then
proves that $M_g$ is flat whenever $M$ is flat.

For the converse retain any injection of $R$-modules
$N'\hookrightarrow N$ and its original $A$-module kernel
$K=\Ker(N'\otimes_R M\to N\otimes_R M)$.
If every $M_{g_\lambda}$ is flat over $R$, localization
and the first tensor comparison identify the localized
map with
$N'\otimes_R M_{g_\lambda}\to N\otimes_R M_{g_\lambda}$,
which is injective. Thus $K_{g_\lambda}=0$ for each $\lambda$.
At a maximal ideal $\mathfrak m$ with $M_{\mathfrak m}=0$,
the same comparison gives
$(N'\otimes_R M)_{\mathfrak m}=N'\otimes_R M_{\mathfrak m}=0$,
so $K_{\mathfrak m}=0$. At any other maximal ideal choose
$g_\lambda\notin\mathfrak m$ by the hypothesis.
For every $k\in K$, the equality $K_{g_\lambda}=0$ gives
$g_\lambda^e k=0$ for some $e\geq0$, which may depend
on $k$. This power is outside $\mathfrak m$, so $k/1=0$
in $K_{\mathfrak m}$. Thus $K_{\mathfrak m}=0$ at every
maximal ideal, and zero detection gives $K=0$.
The injection criterion proves flatness over $R$.

The stated finite unit-ideal hypothesis implies this cover:
if every $g_i$ belonged to a maximal ideal, so would their
generated ideal, contrary to $(g_1,\ldots,g_m)=A$.
This proves (5), including both directions.

The stronger covering hypothesis can equivalently be
stated on all primes $\mathfrak q$ with $M_{\mathfrak q}\ne0$.
Indeed choose $x\in M$ with $x/1\ne0$ at such a prime.
Then $\operatorname{Ann}_A(x)\subseteq\mathfrak q$ by
the exact fraction-zero test. At every maximal
$\mathfrak m\supseteq\mathfrak q$, the same inclusion
shows $x/1\ne0$ in $M_{\mathfrak m}$. A cover element
outside $\mathfrak m$ is also outside $\mathfrak q$.
The reverse implication specializes the prime condition
to maximal ideals. No finite generation of $M$ or of the
family is required.

This really weakens the unit-ideal requirement. For a
nonzero ring $R$, take $A=R\times R$ with the diagonal
map from $R$, and an $R$-module $M$ viewed as an $A$-module
through the first projection. The element $g=(1,0)$ acts
as the identity on $M$, so $M_g=M$ by the maps
$m/g^n\mapsto m$ and $m\mapsto m/1$.
The element $1-g=(0,1)$ acts as zero. If a prime contains
$g$, it cannot contain $1-g$, which becomes an invertible
annihilator of the localized module; hence that localization
is zero. Thus $D(g)$ covers every nonzero stalk of $M$,
although $gA=R\times0$ is a proper ideal of $A$.

Finally all assertions include zero cases. If $0$ lies
in a multiplicative set, its localized ring and modules
are zero by the fraction rule; a unital module over that
ring is zero. In particular $g=0$ or a nilpotent $g$ gives
$M_g=0$ and an empty principal open. If $A=0$, then $M=0$
and its prime and maximal conditions are empty. If the
cover family is empty, its hypothesis says all maximal
stalks of $M$ are zero, so $M=0$ by zero detection and the
equivalence is valid. The zero module is flat by the
definition in every case.
\end{proof}

\begin{lemma}
\label{editorial-source-lemma-flat-going-down}
Let $R \to S$ be flat. Let $\mathfrak p \subset \mathfrak p'$
be primes of $R$. Let $\mathfrak q' \subset S$ be a prime of $S$
mapping to $\mathfrak p'$. Then there exists a prime
$\mathfrak q \subset \mathfrak q'$ mapping to $\mathfrak p$.
\end{lemma}

\begin{proof}
Let $\varphi:R\to S$ be the original map and retain the
stated primes. The localization map
$$
\varphi':R_{\mathfrak p'}\longrightarrow S_{\mathfrak q'},
\qquad r/t\longmapsto\varphi(r)/\varphi(t)
$$
is defined because $t\notin\mathfrak p'$ implies
$\varphi(t)\notin\mathfrak q'$. Its maximal-ideal
contraction is $\mathfrak p'R_{\mathfrak p'}$, so it is
local, and Lemma \ref{editorial-source-lemma-flat-localization} makes it
flat. It is faithfully flat by Lemma \ref{editorial-source-lemma-local-flat-ff}.
The ideal $\mathfrak pR_{\mathfrak p'}$ is prime and has
contraction $\mathfrak p$ in $R$ by prime localization.
Lemma \ref{editorial-source-lemma-ff-rings} supplies a prime
$\widetilde{\mathfrak q}\subset S_{\mathfrak q'}$ whose
contraction under $\varphi'$ is that ideal. Define
$$
\mathfrak q=\{s\in S:s/1\in\widetilde{\mathfrak q}\}.
$$
It is prime as an inverse image of a prime. An element of
$S\setminus\mathfrak q'$ becomes a unit in $S_{\mathfrak q'}$
and cannot belong to this inverse image, so
$\mathfrak q\subseteq\mathfrak q'$. Finally
$$
r\in\varphi^{-1}(\mathfrak q)
\Longleftrightarrow
r/1\in\mathfrak pR_{\mathfrak p'}
\Longleftrightarrow r\in\mathfrak p.
$$
This proves the exact required contraction, without
assuming the original flat map is injective.
\end{proof}

\noindent
The property of $R \to S$ described in the lemma is called the
``going down property''. See Definition \ref{algebra-definition-going-up-down}.

\begin{lemma}
\label{editorial-source-lemma-colimit-faithfully-flat}
Let $R$ be a ring. Let $\{S_i, \varphi_{ii'}\}$ be a directed system of
faithfully flat $R$-algebras. Then $S = \colim_i S_i$ is a faithfully flat
$R$-algebra.
\end{lemma}

\begin{proof}
The underlying $R$-module colimit is flat by
Lemma \ref{editorial-source-lemma-colimit-flat}. For every maximal ideal
$\mathfrak m\subset R$, faithful flatness at each stage
makes $S_i/\mathfrak mS_i$ nonzero, by Lemma \ref{editorial-source-lemma-ff}.
The canonical ring comparison is
$$
\colim_i(S_i/\mathfrak mS_i)\longrightarrow S/\mathfrak mS,
\qquad [s_i+\mathfrak mS_i]\longmapsto[s_i]+\mathfrak mS.
$$
It is surjective by representatives. If $[s_i]\in\mathfrak mS$,
write its original expression as a finite sum
$[s_i]=\sum_{\ell=1}^q a_\ell[t_{j_\ell}]$ with
$a_\ell\in\mathfrak m$. Move all terms to a common stage
and then to a further stage where this equality holds.
There the image of $s_i$ belongs to the actual ideal
$\mathfrak mS_k$, so its quotient class is eventually zero.
This proves injectivity and gives the inverse by choosing
any original representative; equality of two choices is
handled by the same kernel argument.

If this colimit were the zero ring, the class of $1$ at
one stage would become zero at some later stage. Every
transition is a unital $R$-algebra map, so the image is
the original $1$ of that later quotient, contradicting
its nonzero ring structure. Hence $S/\mathfrak mS\ne0$
for every maximal ideal, and Lemma \ref{editorial-source-lemma-ff} gives
faithful flatness. This is the specific reason algebra
colimits preserve faithful flatness although arbitrary
module colimits need not: the same element $1$ survives
in every receiving residue quotient. For $R=0$ all
unital $R$-algebras and their modules are zero, and the
faithfully-flat conclusion holds by the definition.
\end{proof}





